\begin{scanpage}{0152}{134}
Calculations of the scattering characteristics were carried out based on the equations derived above. These scattering characteristics are the functions $h(\alpha,\phi)$ and $e(\alpha,\phi)$ determining the fringing field by means of the relationships
\[
\left.\begin{aligned}E_z&=E_{0z}kae(\alpha,\varphi)\sqrt{\frac2{\pi kr}}e^{i(kr+3\pi/4)},\\
H_z&=H_{0z}kah(\alpha,\varphi)\sqrt{\frac2{\pi kr}}e^{i(kr+3\pi/4)}.
\end{aligned}\right\}\tag{23.13}
\]
The calculations were performed for the values $ka=\sqrt{28}$ and $ka=\sqrt{80}$. In Figures 49--62, the following designations were used: 1) the functions h and e correspond to the rigorous theory; 2) the functions $h_0$ and $e_0$ correspond to the field from the uniform part of the current (the physical optics approach); 3) the functions $h_1$ and $e_1$ correspond to the field from the uniform and nonuniform parts of the current, but without consideration of the interaction of the edges; 4) the functions $h_2$ and $e_2$ correspond to the fringing field with consideration of secondary diffraction calculated on the basis of equations (23.13), (23.02) and (23.10). Thus, in accordance with \S~6,
\[
\left.\begin{aligned}e_0&=\cos\alpha\frac{\sin[ka(\sin\alpha-\sin\varphi)]}{ka(\sin\alpha-\sin\varphi)},\\
h_0&=\cos\varphi\frac{\sin[ka(\sin\alpha-\sin\varphi)]}{ka(\sin\alpha-\sin\varphi)}
\end{aligned}\right\}\tag{23.14}
\]
and
\[
\left.\begin{matrix}e_1\\h_1\end{matrix}\right\}=\frac1{2ka}\left\{\frac{\sin[ka(\sin\alpha-\sin\varphi)]}{\sin\frac{\alpha-\varphi}{2}}\pm i\frac{\cos[ka(\sin\alpha-\sin\varphi)]}{\cos\frac{\alpha+\varphi}{2}}\right\},\tag{23.15}
\]
where $|\alpha|<\pi/2$; $|\varphi|<\pi/2$.

The results obtained show that our approximation equations agree satisfactorily with the rigorous theory already when $ka=\sqrt{28}$, although in the given case approximately one and one-half wavelengths are fitted into the width of the strip. In the direction toward the source ($\varphi=-\pi+\alpha$, $0<\alpha<\pi/2$), and also with glancing irradiation of the
\end{scanpage}
